\input{technique-note.tex}
\title{Integration Bee --- Unit 02:\\ Symmetry \& Basic Trig Integrals}
\author{UK Uni Integration Bee prep\\ Official: integration.soc.srcf.net\\ Yang Topics Guide 2nd Ed §§1.3, 1.5}
\date{}
\begin{document}
\maketitle

\section*{Status}
Concept teaching only (easy drills). No past papers.
Tags: \texttt{odd-even}, \texttt{symmetry}, \texttt{basic-trig}, \texttt{yang:1.3}, \texttt{yang:1.5}.

\section{Official anchors}
\begin{itemize}
  \item Syllabus §1 --- Odd and even functions; A-Level/STEP trig integration
        (\texttt{sources/UK\_University\_Integration\_Bee\_Syllabus.pdf}, p.1).
  \item Techniques Guide §1.1 Odd and Even
        (\texttt{sources/UK\_University\_Integration\_Bee\_Techniques\_Guide.pdf}, Contents p.3 / body p.3).
  \item Techniques Guide §1.8 Trigonometric Identities (Contents p.18) --- deeper identities later;
        this unit stays at basic antiderivatives / powers with $u$-sub.
  \item Yang Topics Guide 2nd Ed §§1.3 Symmetry, 1.5 Basic Trig Integrals
        (\texttt{sources/user/Integration\_Bee\_Topics\_Guide\_Yang\_2nd\_Ed.pdf}).
\end{itemize}

\section{Learning goals}
\begin{enumerate}
  \item Spot odd/even integrands on symmetric intervals and skip useless work.
  \item Split a sum into odd + even parts when helpful (Yang 1.3 style).
  \item Evaluate standard basic trig definite integrals ($\sin$, $\cos$, $\sec^2$, $\tan$, double-angle).
\end{enumerate}

\section{Pedagogy}
Worked example $\to$ faded example $\to$ interleaved easy drills. See \texttt{pedagogy/README.md}.

\section{Worked example A --- odd $\Rightarrow$ 0 (\texttt{U02-D01}, \texttt{yang:1.3})}
\[
\int_{-\pi}^{\pi}\sin x\,dx.
\]
$f(x)=\sin x$ is odd and the interval is symmetric about 0, so the integral is $0$
(Guide §1.1 / Yang 1.3; matches the $x^3$ bridge from Unit 01).

\section{Worked example B --- even doubling (\texttt{U02-D02}, \texttt{yang:1.3})}
\[
\int_{-1}^{1}x^2\,dx = 2\int_0^1 x^2\,dx = 2\cdot\frac13 = \frac23.
\]

\section{Faded example --- split odd+even (\texttt{U02-D03}, \texttt{yang:1.3})}
\[
\int_{-1}^{1}(x^3+x^2)\,dx.
\]
Scaffold: $x^3$ odd $\Rightarrow$ contribution $0$; $x^2$ even $\Rightarrow$ $2\int_0^1 x^2\,dx=\frac23$.

\section{Yang-style expand (\texttt{U02-D04}, \texttt{yang:1.3})}
\[
\int_{-1}^{1}(x+1)(x^2+1)\,dx = \int_{-1}^{1}(x^3+x^2+x+1)\,dx.
\]
Odd terms vanish; even part $2\int_0^1(x^2+1)\,dx=\frac83$ (Yang 1.3 Example 1 shape).

\section{Worked trig --- \texttt{U02-D06} (\texttt{yang:1.5})}
\[
\int_0^{\pi/4}\sec^2 x\,dx = \tan x\Big|_0^{\pi/4}=1.
\]

\section{Faded trig --- \texttt{U02-D08} (\texttt{yang:1.5})}
\[
\int_0^{\pi/4}\tan x\,dx = \int_0^{\pi/4}\frac{\sin x}{\cos x}\,dx.
\]
Scaffold: $u=\cos x$, $du=-\sin x\,dx$, limits $1\to\frac{\sqrt2}{2}$:
\[
-\int_1^{\sqrt2/2}u^{-1}\,du = \ln\sqrt2 = \frac{\ln 2}{2}.
\]

\section{Practice}
Attempt \texttt{U02-D01}\ldots\texttt{U02-D10} in \texttt{drills.tex}.
Machine check: \texttt{verify/run.sh}. CLI tags/answers: \texttt{drills.md}, \texttt{verify/drills\_unit2.json}.

\section{Skips (later)}
King's/Queen's (Yang 1.18 / Guide §1.2), intermediate trig (Yang 1.15), Weierstrass, past Round 1 papers.
\end{document}
